\sgaStableTarget{sga6.e09.expose}\section*{\underline{EXPOSÉ IX}\\[3pt] \underline{\underline{SOME COMPUTATIONS OF GROUPS @@M0@@}}}
\begin{center}
\underline{by P. BERTHELOT}
\end{center}

This exposé corresponds to no oral exposé, and will not be used in the rest of the seminar. It contains a number of computations of groups @@M1@@ for various fibrations: vector bundles, projective bundles, torsors under a linear group, etc\ldots{} Throughout this exposé, the base scheme will be supposed noetherian, and the notation @@M2@@ will denote the Grothendieck group of coherent sheaves on a scheme @@M3@@. Recall that if @@M4@@ is noetherian, the homomorphism from the @@M5@@ of the category of coherent @@M6@@-Modules to the @@M7@@ of the category of pseudo-coherent complexes on @@M8@@ is an isomorphism (IV 2.4).

The results presented here are analogous to those obtained for the Chow ring, which can be found in Grothendieck's exposés in the Chevalley Seminar 1958. The question also arises of determining to what extent statements of the same type can be obtained for the graded group associated with @@M9@@ for the filtration defined by the dimension of the support (cf. X). It does not seem that the answer is known, even for Proposition 1.1 (homotopy exact sequence), which is one of the properties most used in the computations that follow.

\sgaStableTarget{sga6.e09.sec.1}\subsection*{1. Vector bundles}

\sgaStableTarget{sga6.e09.prop.1-1}\textbf{Proposition 1.1. (Homotopy exact sequence).} \uline{Let} @@M10@@ \uline{be a noetherian scheme}, @@M11@@ \uline{a closed subscheme of} @@M12@@, @@M13@@ \uline{the open complement of} @@M14@@ \uline{in} @@M15@@, @@M16@@ \uline{the immersion of} @@M17@@ \uline{into} @@M18@@, \uline{and} @@M19@@ \uline{the immersion of} @@M20@@ \uline{into} @@M21@@. \uline{Then the sequence}:
